\begin{abstract}
QSAR models for Ames mutagenicity are accepted as supporting evidence under EU REACH, 
K-REACH, and ICH M7, yet their reported performance varies widely across evaluation 
protocols, making it difficult to assess their true predictive 
reliability.\cite{Honma_2019,Furuhama_2023,Wu_MoleculeNet_2018} This study systematically 
evaluated 225 model configurations---six tabular algorithms paired with four feature 
representations, plus a graph convolutional network with graph features, under three 
preprocessing scenarios---across three genotoxicity datasets: Ames ($n = 13{,}560$), 
in-vitro chromosome aberration ($n = 1{,}619$), and in-vivo micronucleus ($n = 2{,}153$). 
Using a single locked configuration (XGBoost, compact features, raw SMILES) in an evaluation 
cascade of increasing rigor, Ames MCC fell from 0.670 (random-split CV) to 0.624 
(scaffold-split CV) to 0.470 (scaffold-split held-out test) and further to 0.234 and 0.122 
(drug$\to$industrial and industrial$\to$drug leave-domain-out, LDO). The split-method effect 
($\Delta = 0.046$) was nearly an order of magnitude smaller than the data source 
heterogeneity effect ($\Delta = 0.390$): a naive source classifier reached 
MCC~$= 0.608$ without any structural information, reflecting a 19.5-fold prevalence gap 
between drug-sourced (58.3\%) and industrial-sourced (3.0\%) compounds. Threshold 
decomposition showed that prior shift accounts for roughly one-third (30.9--35.3\,\%) of the 
LDO performance gap, leaving a substantial residual that no recalibration can eliminate. In-domain algorithm differences were small (MCC range 0.063), but 
robustness under source shift was uneven: Random Forest maintained MCC~$= 0.325$ in the 
hardest direction where XGBoost collapsed to 0.122---suggesting that algorithm choice 
matters far more under cross-source evaluation than in-domain comparisons indicate. In-vivo 
micronucleus models showed ranking ability (ROC-AUC~$= 0.860$, PR-AUC~$= 0.237$) but 
classification metrics with 95\% CIs crossing zero, reflecting both modeling limits and 
limited test-set positives ($n = 12$). These findings support an evaluation cascade with 
source-awareness as a useful addition to current reporting practices for mutagenicity QSAR.
\end{abstract}
